\documentclass{homework}
% \usepackage{graphicx} % Required for inserting images
\usepackage{amsmath}
\usepackage{gensymb}
\usepackage{subfig}
\usepackage{float}

\class{ECEN 5427}
\title{Project 2 Proposal: The Impact of Battery Degradation on Power System Operations}
\author{Sean Jennings}
\date{March 11, 2025}

\begin{document}
\maketitle

\section{Introduction}

For this project, I plan to quantify the system-wide benefits of grid-scale Lithium Ion Battery (LiB) storage, with a focus on battery lifetime and degradation modeling.
LiB degradation is a complex electrochemical phenomenon, but as LiBs have developed into a commercially mature technology, a vast array of research has uncovered
many useful insights (see [] for an overview). I am interested in exploring the implications of degradation to questions of systems-level operations.

Many techno-economic analyses of battery energy storage systems (bESSs) amortize capital costs over a fixed lifetime of typically 10-15 years for LiB packs. However,
many studies show that at least at the smallest scale of assembly topologies, battery cells age both as a function of
\begin{enumerate}
    \item cycle aging: that is, the conditions under which the cell are cycled, especially
    \begin{itemize}
        \item depth of discharge (DoD),
        \item max (dis)charge rates (also known as C-rates); and
    \end{itemize}
    \item calendar aging: i.e. how the battery deteriorates over time, which also depends on
    \begin{itemize}
        \item the state of charge (SoC) at which the battery is stored
    \end{itemize}
\end{enumerate} 

It happens that these operating conditions are (more or less directly) inputs to system-level problems of economic dispatch (ED) and unit commitment (UC), which suggests that
it makes sense to amortize over battery usage, rather than over time, and treat usage as an operational cost, akin to the fuel costs for thermal generators.

Taking a control perspective for a moment, we will use the term state of health (SoH) to refer some unknown representation of the internal electrochemical
condition of the battery. The SoH is observable in several ways:
\begin{enumerate}
    \item capacity fade: the amount of charge which can be held diminishes with declining SoH,
    \item power fade: the maximum discharge and charge rates decay with declining SoH,
    \item open Circuit Voltage (OCV) curves changes: OCV is the cell voltage as a function of SoC (when allowed to (dis)charge sufficiently slowly) and provides information about internal electrochemical mechanism
    \item electrochemical Impedance Spectroscopy (EIS): also correlates closely with battery state of health
\end{enumerate}

I plan to focus mostly on capacity fade, but for now I will note that the OCV curves and EIS could be very useful in online applications. By monitoring SoH continuously, bESS facilities and
their owners will be able to improve their degradation and lifetime prediction over time, and ideally this should feedback into capacity expansion models. But I will consider
closed-loop control and state estimation problems as out of scope for this project.

Unfortunately, the Energy Storage "External" which SAint provides is a bit limited in its reprentation. There are a number of UC-related parameters (mininum up/downtime, etc.) which 
are likely not very relevant since LiB switching speeds are fast (see LiB's use frequency response applications). There is a field for "variable operation and maintenance cost", but it only
seems to support a fixed value, rather than a power-dependent curve. Instead, I am working on model which I believe will more accurately represent grid-scale bESS facilities, by including
the coupling between several physical domains:
\begin{enumerate}
    \item a simple electrochemical model of degradation,
    \item a heat transfer model to determine the amount of efficiency lost to temperature control systems as a function of ambient temperature
    \item the electrical domain which connects the bESS to the grid and our DC optimal power flow problem (DCOPF)
\end{enumerate}

\section{Methodology}

I plan to modify the 37 bus network provided by the encoord to seek answers to the following questions:

\begin{letterenum}
\item Q1: how does a higher-order degradation models affect the solutions of the ED and UC problems?
\begin{itemize}
    \item Does it predict a lower/higher overall cost?
    \item 
\end{itemize} 
\item Q2: Potential for BESS to alleviate transmission constraints
\end{letterenum}

\section{Problem Formulation}





Temperature is also a major factor in determining degradation rates. 
Unfortunately, SAint does not 

\begin{enumerate}
\item Research Questions
\item Methodology
    \begin{letterenum}
        \item Production Cost Modeling:
            \begin{itemize}
            \item Unfortunately, SAint doesn't appear to have the capabilities to explore the first question: fixed depth of discharge (aka min/max SoC)
            \item Add a mock-up 2D fuel cost curve that exemplifies the issue
            \item Breakdown into multiple energy storage objects to effect a piecewise linear model (Well, I don't actually think this works since SoC constraints are interdependent...)
            \item Possibly, need to write custom PCM optimization model to test this properly
            \end{itemize}
        \item Data Sources:
        \begin{itemize}
            \item BESS Pricing Inputs: breakdown capital costs by component (Power Electronics vs. Battery packs vs. others?)
            \item Battery pack degradation (Lithium ion cells) -> operating cost
            \item Renewable Energy Profiles: use New England-based solar/wind resources?
        \end{itemize}
        \item Test Network:
        \begin{itemize}
            \item There's enough complexity in the initial part of the 
        \end{itemize}
        \item Scenarios
        \begin{enumerate}
        \item Q1: Consider more of a scenario matrix:
            \begin{itemize}
            \item Base Case -- high renewable penetration with simple (zero-order) O\&M cost model
            \item Zero-order vs first-order modeling
            \item Renewable penetration (low, mid, high)
            \item Energy Storage sizing
            \item Cost sensitivities -- Lithium prices (low, mid, high)
            \end{itemize}
        \item Q2: Transmission constraints
            \begin{itemize}
            \item Base Case -- Moderate renewable penetration, transmission-constrained system
            \item Energy Storage location --> colocated with renewable generation, baseload generation, demand, or near transmission constraint
            \end{itemize}
        \end{enumerate}
    \end{letterenum}
\end{enumerate}

For this project, I am interested in exploring and understanding the operational impacts of integrating (eletrochemical-based) battery energy storage systems (BESS's) into power systems with high penetrations
of renewable energy resources. 

\section{Research Questions}

My first question pertains to the degradatory nature of Lithium Ion-based BESS's. Batteries experience irreservible electrochemical changes which cause their storage capacity and power generation
capabilities to reduce over time.  Typically, when assessing the upfront capital investment of a BESS, the power electronic components are specc'ed out for 20 year lifetimes, while the
battery packs anticipate a 10 year lifetime. While this approximation may be satisfactory for capacity expansion-like scenarios, where uncertainty is high, the true picture is quite
a bit more complex, and power system operators (and planners) may benefit by considering a more nuanced picture of BESS lifetime degradation. Specifically, battery degradation is typically
broken down into calendar aging and cycle aging. Models which assume a fixed 10 year lifetime implicitly neglect the cycle aging component. The cycle aging component is complex, since
it depends on a number of factors which depend both on the BESS design as well as the operating conditions (i.e. power output and the depth of discharge). A power system which accurately
models the cycling-related degradation dependencies should be able to achieve a dispatch with lower system-wide costs. 

Essentially, cycle aging can be considered as the marginal cost of a BESS, analogous to the fuel costs of a thermal generator.

This leads to our first research question:


{\bf Question 1:} How does 



The first focuses on the near-term potential
of providing transmission congestion relief. 

This project aims to develop a production cost model (PCM) in SAInt to analyze the economic and operational impacts of energy storage degradation on power system operations.
A comparative analysis will be conducted between a simplified first-order approximation of battery degradation and a detailed empirical model accounting for both cycling and calendar aging.

\section{Research Questions}
Research Questions:

    How close is a first-order approximation of battery lifespan to reality?
        Does incorporating detailed degradation effects lead to an over- or under-estimation of total system costs compared to a simplified model?
    How does economic dispatch differ between the two models?
        Does the more detailed degradation model shift optimal storage dispatch strategies?
        How does it impact renewable integration and system-level costs?

\section{Methodology}

    Test System Selection
        Start with encoord’s test system (20+ buses with renewables).
        Integrate real-world data from EIA’s AEO and EIPC to make the model more representative.

    Production Cost Modeling in SAInt
        Formulate unit commitment + economic dispatch with DC power flow.
        Simulate hourly operations for one year under different scenarios.
        Use a simplified first-order degradation model for battery lifespan.

    Extending the Model Beyond SAInt (If Possible)
        If SAInt doesn’t support a detailed degradation model, consider building an external PCM in CVXpy for comparison.
        Use empirical battery degradation models to estimate real-world operational impacts.

    Scenario Design (At least 5 per research question)
        Baseline: Standard SAInt PCM with first-order degradation.
        Detailed Degradation Model: PCM with empirical degradation effects.
        High Storage Penetration: More storage to test sensitivity.
        High Renewable Penetration: More wind/solar to observe dispatch changes.
        Different Battery Chemistries: If time allows, compare Li-ion to other storage tech.




\end{document}
